\section{Hipervisors}

Um hipervisor, também conhecido como \gls{vmm}, é uma plataforma que cria e executa máquinas virtuais via virtualização.
Segundo \citeonline{PopekGeraldGoldberg} um \gls{vmm} possuí três características essenciais: primeiro, devem fornecer um ambiente idêntico à máquina original
para programas; segundo, os programas executados nesse ambiente, na pior das hipóteses, devem apresentar no máximo pequenas reduções
de velocidade que urgem da camada de virtualização; e por último, os \glspl{vmm} devem deter o contole total sobre os recursos do sistema.

Na visão de \cite{Goldberg}, fundamentalmente, existem dois tipos de hipervisores: 1 e 2.

\subsection{Tipo 1 - Bare Metal}

Hipervisores de tipo 1, também conhecidos como bare metal, são aqueles que executam diretamente sobre o hardware físico, sem a necessidade de um
sistema operacional hospedeiro intermediário.
Por operar diretamente sobre o hardware, esse tipo de hipervisor geralmente apresenta melhor desempenho e menor sobrecarga.

\begin{figure}[H]
    \centering
    \begin{tikzpicture}[
        pequeno/.style={
            draw,
            thick,
            fill=blue!5,
            minimum width=4cm,
            minimum height=1.2cm,
            outer sep=0pt
        }
    ]
        \node[pequeno] (apps-so1-1) {Aplicação 1};
        \node[pequeno, right=0.2cm of apps-so1-1] (apps-so1-2) {Aplicação 2};
        \node[pequeno, right=0.2cm of apps-so1-2] (apps-so2-1) {Aplicação 1};

        \node[fit=(apps-so1-1)(apps-so1-2), inner sep=0pt] (apps-so1) {};
        \node[fit=(apps-so2-1), inner sep=0pt] (apps-so2) {};

        \node[
            draw,
            thick,
            fill=blue!10,
            minimum width=8.2cm,
            minimum height=1.5cm,
            below=0.2cm of apps-so1
		] (SO1) {\Gls{so} 1};

		\node[
            draw,
            thick,
            fill=blue!10,
            minimum width=4cm,
            minimum height=1.5cm,
            below=0.2cm of apps-so2
		] (SO2) {\Gls{so} 2};

		\node[fit=(SO2)(SO1), inner sep=0pt] (SOs) {};

		\node[
            draw,
            thick,
            fill=blue!15,
            minimum width=12.4cm,
            minimum height=1.5cm,
            below=0.2cm of SOs
		] (HOS) {Hipervisor};

		\node[
            draw,
            thick,
            fill=blue!20,
            minimum width=12.4cm,
            minimum height=1.5cm,
            below=0.2cm of HOS
		] () {Hardware};

    \end{tikzpicture}
    \caption{Hipervisor Tipo 1.}
    \label{fig:hypervisor-type-1}
\end{figure}

\subsection{Tipo 2}

Diferentemente da solução bare metal, um hipervisor do tipo 2 é executado como um processo dentro de um \gls{so} hospedeiro.
Porém, apesar de garantir os mesmos recursos do tipo 1, sua arquitetura impõe uma camada intermediária a mais de software, o que restringe a
eficiência no acesso aos componentes físicos, gerando uma latência adicional e um menor desempenho.

\begin{figure}[H]
    \centering
    \begin{tikzpicture}[
        pequeno/.style={
            draw,
            thick,
            fill=blue!5,
            minimum width=4cm,
            minimum height=1.2cm,
            outer sep=0pt
        }
    ]
        \node[pequeno] (apps-so1-1) {Aplicação 1};
        \node[pequeno, right=0.2cm of apps-so1-1] (apps-so1-2) {Aplicação 2};
        \node[pequeno, right=0.2cm of apps-so1-2] (apps-so2-1) {Aplicação 1};

        \node[fit=(apps-so1-1)(apps-so1-2), inner sep=0pt] (apps-so1) {};
        \node[fit=(apps-so2-1), inner sep=0pt] (apps-so2) {};

        \node[
            draw,
            thick,
            fill=blue!10,
            minimum width=8.2cm,
            minimum height=1.5cm,
            below=0.2cm of apps-so1
		] (SO1) {\Gls{so} 1};

		\node[
            draw,
            thick,
            fill=blue!10,
            minimum width=4cm,
            minimum height=1.5cm,
            below=0.2cm of apps-so2
		] (SO2) {\Gls{so} 2};

		\node[fit=(SO2)(SO1), inner sep=0pt] (SOs) {};

		\node[
            draw,
            thick,
            fill=blue!15,
            minimum width=12.4cm,
            minimum height=1.5cm,
            below=0.2cm of SOs
		] (HOS) {Hipervisor};

		\node[
            draw,
            thick,
            fill=blue!20,
            minimum width=12.4cm,
            minimum height=1.5cm,
            below=0.2cm of HOS
		] (OS) {\Gls{so}};

		\node[
            draw,
            thick,
            fill=blue!25,
            minimum width=12.4cm,
            minimum height=1.5cm,
            below=0.2cm of OS
		] () {Hardware};

    \end{tikzpicture}
    \caption{Hipervisor Tipo 2.}
    \label{fig:hypervisor-type-2}
\end{figure}
